% !TeX root = ../../packet-and-principal-models.tex
\section{The complete packet: arithmetic and geometry}
The proof separates the finite arithmetic domain from curve membership.
The supplementary order calculations feed into the following full census;
its proof includes noninvertible modules over the originally generated order.
Historical partial exclusion totals are kept in the supplement and play no
role as substitutes for the final two-target cover.

\subsection{The complete marked arithmetic census}\label{subsec:arithmeticcensus}

Throughout this subsection the extension $M/L$, its $L$-action and the
ordered six real embeddings are fixed.  Write $z=\zeta_{42}$ and
\[
 \alpha_k=z^k(1+z+z^4)=A_k+B_kz,\qquad
 c_k=2A_k+tB_k,\qquad 0\leq k<42.
\]
We count integral matrices of trace $c_k$ and determinant $\pi$, up to
$\SL_2(\OL)$ conjugacy, and their upper fixed tuples in
$\SL_2(\OL)\backslash\HH^6$.  Thus the universe is the norm-$\pi$ packet
in this marked CM field, rather than a union over other CM fields.
For an ideal $\mathfrak f\subseteq\OL$, put
$S_{\mathfrak f}=\OL+\mathfrak f\OM$; in particular
$S_{\OL}=\OM$.

\begin{lemma}[Point fibers and the intrinsic multiplier order]
\label{lem:packetpointfibers}
For any fixed trace $c_k$, the map from its integral matrix classes to
quotient upper fixed points has exactly two classes in each fiber,
represented by $T$ and $c_kI-T$.  The relative conductor ideal of the
multiplier order is an invariant of the quotient point.  In particular,
different embedded ideals in Proposition~\externalref{supplement}{prop:alltraceorders}
give disjoint proper-multiplier point populations.
\end{lemma}
\begin{proof}
The two-root argument in Theorem~\externalref{supplement}{thm:d7ambient} applies to every
$c_k$: the characteristic polynomial is irreducible, since its
discriminant $B_k^2(t^2-4)$ is negative at every real embedding.
If two actions fix the same upper tuple, they preserve both conjugate
complex lines and commute.  Injectivity of a real embedding makes this
commutation an identity over $L$.  They therefore belong to the same
quadratic centralizer, where their characteristic polynomial has only
its two roots.  An anti-intertwiner has determinant $-N_{M/L}(a)$,
totally negative, so the two root classes are distinct under $\SL_2$.

The $L$-algebra preserving those lines is intrinsic to the tuple.
Its intersection with $M_2(\OL)$ is the multiplier ring of the
corresponding lattice.  An integral determinant-one change of basis
conjugates both the algebra and this order.  The two identifications of
that algebra with $M$ fixing $L$ differ only by complex conjugation,
which preserves every $S_{\mathfrak f}$.  Hence the exact ideal
$\mathfrak f$, not merely its norm, is invariant.  Extra endomorphisms
of an unmarked abelian variety do not enlarge this two-by-two
$L$-linear centralizer, and automorphisms permuting the marked real
embeddings are not quotient identifications.
\end{proof}

\begin{proposition}[The $144$ selected $d=3$ ambient points]
\label{prop:d3ambientpoints}
The $1280$ matrix classes of Theorem~\externalref{supplement}{thm:d3classes} define exactly
$144$ distinct quotient points.  Their decomposition and the number
of those matrix classes incident to each point are
\[
\begin{array}{c|r|r}
\text{Multiplier conductor}&\text{Points}&\text{Incidences per point}\\\hline
\OL&16&24\\
(7,t-1)&32&16\\
(7,t-1)^2&32&4\\
(43,t-20)&32&4\\
(43,t-31)&32&4.
\end{array}
\]
Their intersection with the $112$-point population of
Theorem~\externalref{supplement}{thm:d7ambient} consists of exactly the sixteen maximal
points; the union contains $240$ distinct ambient points.
\end{proposition}
\begin{proof}
Theorem~\externalref{supplement}{thm:d3classes} supplies the complete module and orientation
lists.  Choose the common basis $(1,gz)$ with
\[
 g=1,\ B_0,\ B_{12},\ B_4,\ B_8
\]
for the five displayed multipliers.  Let $\delta_m$ be the binary
product selected by mask $m$ from
\[
 (\xi_5,\xi_{11},\xi_{13},\xi_{17},-1,u),\qquad u=4-t^2.
\]
Use masks $0,\ldots,31$ for maximal multiplier and $0,\ldots,63$
for a proper multiplier.  Multiplication by any allowed $\alpha_k$
in orientation $\delta=\delta_m$ is
\[
 T_{k,g,\delta}=
 \begin{pmatrix}A_k&-\delta gB_k\\
 B_k/(\delta g)&A_k+tB_k\end{pmatrix}.
\]
Putting $q=g\delta$, the upper tuple is specified exactly by
\begin{equation}\label{eq:principalpacketpoint}
 2x=-qt,\qquad 4h^2=q^2(4-t^2),\qquad |w|^2=q^2,
 \qquad h>0.
\end{equation}
It is independent of the trace label.  The matrix
\[
 U_q=\begin{pmatrix}1&qt\\0&1\end{pmatrix}\in\SL_2(\OL)
\]
sends this tuple to the one for $-q$ and satisfies
$U_q(c_kI-T_{k,g,\delta})U_q^{-1}=T_{k,g,-\delta}$.
Thus root pairing toggles bit $16$.  Canonical masks are $0,\ldots,15$
for $\OM$ and $0,\ldots,15,32,\ldots,47$ for every proper multiplier.
Lemma~\ref{lem:packetpointfibers} proves that no further quotient
identifications occur.  Literal equality of the chosen tuples is a
separate assertion: their real parts force equality of $q$, whose ideal
$(q)=(g)$ separates the five strata.

We also account for the basis orientations in the old $1280$-matrix
list.  A regular basis $(1,\alpha_k)$ changes to $(1,gz)$ by
$Q=\left(\begin{smallmatrix}1&A_k\\0&e\end{smallmatrix}\right)$,
where $e=B_k/g$; for the intermediate basis use
$Q=\operatorname{diag}(1,(t-1)/B_0)$.  The nontrivial unit-square
identities $e=\delta_a v^2$ are
\[
\begin{array}{c|r|l}
 e&a&v\\\hline
 B_{15}/B_0&33&1+5t-5t^3+t^5\\
 B_{18}/B_0&39&3+13t-t^2-11t^3+2t^5\\
 (t-1)/B_0&7&1+18t+3t^2-16t^3-t^4+3t^5.
\end{array}
\]
Each displayed $v$ is a unit.  For regular labels $0,4,8,12$ use
$e=v=1,a=0$; the opposite trace label negates $e$ and toggles bit $16$
of $a$.  The intermediate basis is the same at labels $12$ and $33$.
If $n=m\mathbin{\oplus}a$ and
$h_0=\prod_{i:a_i=1,\,m_i=0}s_i$ for the six listed generators $s_i$,
then $\delta_n\delta_a/\delta_m=h_0^2$.  Consequently
\[
 P=(vh_0)^{-1}D_{\delta_n}QD_{\delta_m}^{-1},
 \qquad D_\delta=\operatorname{diag}(\delta,1),
\]
is integral of determinant one and transports the old action to its
common-basis action.  If needed, compose with $U_{g\delta_n}$ to
select the canonical root partner.  The paired exact implementations
in Supplement, Section~\externalref{supplement}{app:computer} check every unit identity and all
$1280$ matrix and upper-tuple transports.

There are twelve trace labels at maximal multiplier, eight at
multiplier $(7,t-1)$, and two at each of the other three proper
multipliers.  The two-root fibers give the displayed incidences and
$16\cdot24+32\cdot16+96\cdot4=1280$.
The intermediate conductor-seven stratum is counted once.
The proper $d=3$ conductors differ from the two proper $d=7$
conductors, whereas their common maximal tuples coincide.
Lemma~\ref{lem:packetpointfibers} gives the asserted intersection
and $144+112-16=240$.
\end{proof}

\begin{theorem}[The full marked norm-$\pi$ arithmetic census]
\label{thm:fullpacketarithmetic}
For all $42$ traces in Proposition~\externalref{supplement}{prop:alltraceorders}, the
complete integral matrix universe has $3520$ trace-labelled
$\SL_2(\OL)$ classes.  The exact counts, including every module
overorder and determinant orientation, are
\[
\begin{array}{c|r|r|r}
\text{Generated-order conductor}&\text{Traces}&
 \text{Classes per trace}&\text{Subtotal}\\\hline
\OL&16&32&512\\
(7,t-1)&12&96&1152\\
(7,t-1)^2&2&160&320\\
(41,t-12)&2&96&192\\
(41,t-22)&2&96&192\\
(43,t-20)&2&96&192\\
(43,t-31)&2&96&192\\
(83,t-36)&2&160&320\\
(127,t-94)&2&224&448.
\end{array}
\]
Their upper fixed tuples define exactly $368$ distinct ambient
quotient points, with the following intrinsic multiplier decomposition:
\[
\begin{array}{c|r|r|r}
\text{Multiplier conductor}&|\operatorname{Pic}(S_{\mathfrak f})|&
 \text{Points}&\text{Incidences per point}\\\hline
\OL&1&16&84\\
(7,t-1)&1&32&28\\
(7,t-1)^2&1&32&4\\
(41,t-12)&1&32&4\\
(41,t-22)&1&32&4\\
(43,t-20)&1&32&4\\
(43,t-31)&1&32&4\\
(83,t-36)&2&64&4\\
(127,t-94)&3&96&4.
\end{array}
\]
In particular the two nonprincipal classes at conductor $(127,t-94)$
are exchanged by conjugation and together contribute $64$ points.
These are arithmetic ambient points; the theorem makes no new
curve-membership or geometric CM-type assertion.
\end{theorem}
\begin{proof}
The existing local calculations in Theorems~\externalref{supplement}{thm:d7nonmax}
and~\externalref{supplement}{thm:d3classes} leave two embedded orders to consider.
Put
\[
 \mathfrak p'_{41}=(41,t-22),\quad
 \mathfrak p_{127}=(127,t-94),\quad \gamma=B_{10}.
\]
Proposition~\externalref{supplement}{prop:alltraceorders} gives
$(B_5)=\mathfrak p'_{41}$ and $(\gamma)=\mathfrak p_{127}$.
For either prime $\mathfrak p=(p,t-r)$, compute in
\[
 A_{\mathfrak p}=\mathbf F_p[Z]/(Z^2-rZ+1),\qquad
 Q_{\mathfrak p}=A_{\mathfrak p}^{\times}/\mathbf F_p^{\times}.
\]
Exact multiplication of projective pairs $[a:b]$, with
$a^2+rab+b^2\ne0$, gives
\[
\begin{array}{c|l|r|c|r}
\mathfrak p&A_{\mathfrak p}&|Q_{\mathfrak p}|&
 [\varepsilon]&\operatorname{ord}[\varepsilon]\\\hline
\mathfrak p'_{41}&\mathbf F_{41^2}&42&[1:30]&42\\
\mathfrak p_{127}&\mathbf F_{127}\times\mathbf F_{127}&126&[1:80]&42.
\end{array}
\]
Here $\varepsilon=1-z^2$ and $[z]=[\varepsilon]^{22}$ in both
quotients.  At $127$ the roots are $40,54$; the ordered ratio of the
two components sends $\varepsilon$ to $51$ and $1+9z$ to $7$.
For example,
\[
\begin{gathered}
 51^6=32,\quad51^{14}=107,\quad51^{21}=-1,\quad51^{42}=1,\\
 7^{18}=64,\quad7^{42}=107,\quad7^{63}=-1,\quad7^{126}=1
 \pmod{127}
\end{gathered}
\]
certify their respective orders $42$ and $126$.
The complete unit generation $E_M=\langle\varepsilon\rangle\mu_{42}E_L$
proved in Theorem~\externalref{supplement}{thm:d7nonmax} shows that no omitted global
unit enlarges the order-$42$ image.  The conductor exact sequence
and $h(M)=1$ give
\[
 \operatorname{Pic}(S_{\mathfrak p'_{41}})=1,
 \qquad \operatorname{Pic}(S_{\mathfrak p_{127}})\simeq\mathbf Z/3\mathbf Z.
\]
For every proper order here, the unit condition is
$n+22j\equiv0\pmod d$ on $\varepsilon^n z^j v$, $v\in E_L$,
where the image order $d$ is $6$ or $42$.  At $127$ it is $42$,
not $|Q_{\mathfrak p_{127}}|=126$.  The exponent $n$ is therefore
even, so $N(S_{\mathfrak f}^{\times})=E_L^2$; conversely all real
unit squares occur.  For maximal multiplier the image is $E_L^+$.

For completeness of the modules at a prime conductor, normalize a
full local rank-one lattice $I$ so its maximal-order span is $B$.
In the inert case a vector of minimal valuation is a unit of $B$;
in the split case the residue subspace surjects to each coordinate
and contains a vector nonzero in both.  Divide by such a unit.
Then $D+\varpi B\subseteq I\subseteq B$, whose quotient has length
one over the local DVR $D$.  Hence there are exactly the proper
and maximal local types.  Globally this gives
\[
 \operatorname{ICM}(S_{\mathfrak p})=
 \operatorname{Pic}(S_{\mathfrak p})\sqcup\operatorname{Pic}(\OM).
\]
The square-conductor proof in Theorem~\externalref{supplement}{thm:d3classes} adds
exactly the intermediate order at seven.  Thus noninvertible modules
over the originally generated order are included.  Since $h(L)=1$,
each lattice has an integral rank-two basis.  The determinant quotient
$E_L/N(S^{\times})$ supplies $32$ classes for maximal multiplier
and $64$ for each proper Picard class.  These facts prove the first table.

To make the new nonprincipal classes and their root fibers explicit,
use the three distinct conductor-unit cosets
$1,1+9z,1+100z$ at $127$ and the ideals
\[
 J_s=\gamma\OM+(1+sz)\OL,
 \qquad (1+sz,\gamma z)\text{ a basis},\qquad s=0,9,100.
\]
Their residues are units at the conductor; locally they are multiples
of the proper order, and away from it they are maximal.  Thus they
are invertible ideals of the proper multiplier order.  The conductor
sequence identifies their three classes and proves they exhaust it.
The action of $\alpha_k=A_k+B_kz$ is
\[
 T_{k,s}=\begin{pmatrix}
 A_k-sB_k&-\gamma B_k\\
 (B_k/\gamma)(1+st+s^2)&A_k+(t+s)B_k
 \end{pmatrix}.
\]
Its entries are integral and its characteristic polynomial is the
required one; conjugating by $D_{\delta_m}$ supplies all orientations.

Conjugation is inversion on $Q_{\mathfrak p_{127}}$, and
$\overline{1+9z}=85(1+100z)$ modulo the conductor.
Consequently $\overline{J_9}=J_{100}$ and conversely.
The conjugation matrix from the basis of $J_9$ to that of $J_{100}$ is
\[
 C_{100,9}=\begin{pmatrix}
 1+9t&\gamma t\\
 -(109+900t)/\gamma&-1-100t
 \end{pmatrix},\qquad \det C_{100,9}=-1.
\]
It is integral, since $109+900\cdot94\equiv0\pmod{127}$ and
$(\gamma)=(127,t-94)$.  Therefore
$D_{\delta_{m\oplus16}}C_{100,9}D_{\delta_m}^{-1}$ has determinant
one and transports the action on $J_9$ to the conjugate-root action
on $J_{100}$.  The root pairs are
$(J_9,m)\leftrightarrow(J_{100},m\oplus16)$.
All $64$ masks on $J_9$ represent the nonprincipal point population;
the principal class contributes $32$ further points, giving $96$.

For clarity, the nonprincipal conductor-$83$ module may be represented
by $(1+3z,B_2z)$.  In its conductor quotient
$[1+3z]^2=[\varepsilon]^{29}$, and
$\eta=\varepsilon z^{28}$ has norm $u$.
Multiplication by $\eta$ followed by conjugation preserves this ideal
and has determinant $-u$.  Thus its root pairing is
$m\mapsto m\oplus48$.  After twisting, multiply the bridge by the
scalar $u^{-1}$ when bit $32$ of $m$ is unset, and by $1$ otherwise;
the determinant is then one.  Principal modules use $m\mapsto m\oplus16$.
The two portable implementations check all $736$ integral root bridges
against every permitted trace action.

Finally, in any basis $(1+sz,gz)$ put
$q=g\delta$ and $n=1+st+s^2=N_{M/L}(1+sz)>0$.
The exact tuple satisfies
\[
 2x=-\frac{q(t+2s)}n,\qquad
 4h^2=\frac{q^2(4-t^2)}{n^2},\qquad |w|^2=\frac{q^2}n.
\]
It is independent of the allowed trace label.  Together with
Lemma~\ref{lem:packetpointfibers}, this identifies repeated overorders
and proves both the point counts and separation of distinct conductor
strata in the second table.  The fourteen traces acting in the
conductor-seven order give $28$ incidences per point; all $42$ act in
the maximal order.  Each other proper stratum occurs for two traces.
Thus the independent incidence check is
\[
 16\cdot84+32\cdot28+320\cdot4=3520,
 \qquad 16+6\cdot32+64+96=368.
\]
The exact manifests also compare all chosen tuple coefficients, but
quotient distinctness follows from the structural arguments above.
\end{proof}

\begin{remark}[Geometry and representative labels]
The $144$-point and $112$-point populations account for $240$ points of
the full census.  The two further conductor strata contribute
$32+96=128$ points, giving $368$.  Equal rational conductor norms do
not identify the two primes over $41$ or the two over $43$.
The common-basis labels used here need not equal older companion or
ideal-lattice masks, especially at nonprincipal conductor $83$.
An individual geometric decision requires an exact representative
transport; a proved whole-stratum conclusion transfers through the
equality of complete intrinsic populations.  This arithmetic census
adds no geometric CM type, polarization or Hodge conclusion.
\end{remark}


\subsection{From two regions to the full curve}
\begin{lemma}[Two simultaneous target regions]\label{lem:twochambers}
Let $Q$ be the labelled quadrilateral chamber in McMullen's source
half-plane, let $r$ be reflection in its $AB$ side, and let $Q_i$ be
its six normalized coordinate quadrilaterals. Write $r_i$ for reflection
in the labelled $A_iB_i$ side and put
\[
 P_0=\prod_{i=1}^6 Q_i,\qquad P_1=\prod_{i=1}^6 r_i(Q_i).
\]
If a tuple $p\in\HH^6$ represents a point of the normalized McMullen
curve, some $H\in\SL_2(\OL)$ sends $p$ into $P_0\cup P_1$.
Consequently, complete empty searches for both target regions exclude
$p$ from the curve. Each search may use its own explicitly verified
$\SL_2(\OL)$ translate of $p$.
\end{lemma}
\begin{proof}
Let $W$ be the quadrilateral reflection group and let $\Delta_0$ be
its orientation-preserving subgroup. McMullen's construction identifies
the holomorphic quadrilateral group with this index-two subgroup
\cite[\S4]{McMullen2026}. Thus $Q\cup rQ$ is a fundamental region
for $\Delta_0$, with boundary identifications allowed. The normalized
modular embedding $F$ is obtained from the labelled triangle maps by
Schwarz continuation. The same continuation across the $AB$ side gives
$F_i(rz)=r_i(F_i(z))$ in every coordinate
\cite[\S6, Theorem~6.3 and the proof of Theorem~1.9]{McMullen2026}.
It follows that $F(Q)\subset P_0$ and $F(rQ)\subset P_1$.
The region choice is simultaneous in all six coordinates; no independent
coordinate reflections are used.

Write $p=KF(z)$ with $K\in\SL_2(\OL)$. Choose
$\eta\in\Delta_0$ such that $\eta z\in Q\cup rQ$.
Equivariance gives
$\rho(\eta)K^{-1}p=F(\eta z)\in P_0\cup P_1$,
where the normalized monodromy lift $\rho(\eta)$ is integral of
determinant one. If search $j$ uses $p_j=G_jp$ with
$G_j\in\SL_2(\OL)$, then $HG_j^{-1}$ carries $p_j$ to the same
target. Both empty exhaustive searches contradict this necessity.
This proof makes no assertion that any $r_i$ normalizes the arithmetic
group or permutes the packet.
\end{proof}

\begin{theorem}[Complete fixed marked geometric packet]\label{thm:fullpacketgeometry}
Assume the packet execution hypothesis stated in the introduction.
Let $\mathcal S$ be the $368$ ambient quotient points of
Theorem~\ref{thm:fullpacketarithmetic}, with its fixed $M/L$, norm
generator, $L$-action and ordered embeddings. Let $\mathcal X$ be the
normalized McMullen curve image in $\SL_2(\OL)\backslash\HH^6$.
Then
\[
 \mathcal S\cap\mathcal X=\{A,B,C\}.
\]
Thus exactly $365$ points are excluded, including every point of proper
multiplier order. The three labels in the arithmetic certificate are
$\texttt{OM:pic0:07}$, $\texttt{OM:pic0:15}$ and
$\texttt{OM:pic0:00}$ for $A,B,C$, respectively.
\end{theorem}
\begin{proof}
The arithmetic census fixes the complete domain, including all module
classes, determinant sectors, intermediate overorders and root fibers.
The positive vertex realizations in Theorem~\ref{thm:A} and
Theorem~\externalref{supplement}{thm:d7ambient} give three distinct points in the intersection.
It remains to exclude each of the other $365$ points.

For the original target $P_0$, the retained legacy calculations cover
$116$ distinct points. The original-target portable certificate covers
the disjoint complement of $249$ nonvertex labels. Their union is exactly
$\mathcal S\setminus\{A,B,C\}$; the thirteen maximal exclusions are already
inside the $116$ and are not counted twice. The original-target coverage
uses the complete numerical searches, not the full-curve conclusions
of the historical subpacket propositions. At nonprincipal conductor
$83$, its bridge is equality of the entire $32$-point population, not
an unproved permutation of the old and new mask labels.

The reflected certificate reconstructs all $365$ nonvertex labels and
tests the simultaneous target $P_1$. In each target the source transport
is separately verified integral of determinant one. Fresh exact target
vertices and outward rational height/modulus bounds give finite
positive trace-form ellipsoids for all possible entries. Inclusive
lattice recursion, the affine determinant equation, both simultaneous
matrix signs and the separate $c=0$ branch exhaust the candidate set.
Closed side predicates retain boundary equality. These arguments apply
to nonprincipal modules and all orientations, without a maximal-order
or mask-zero hypothesis.

The primary quotient-field path and the separate Sturm path agree on
the complete preparations, entries, rows, candidates and classifications.
Every survivor set is empty. The retained reports bind each plan and
result to its source, target, runtime and full row journal. The
whole-executable correspondence and actual portable controls justify
the source relocation; controls are not a sampling proof for untested
labels. The certificates and historical execution assumptions are
specified in Supplement, Section~\externalref{supplement}{app:computer}.

If a remaining point lay on the curve, Lemma~\ref{lem:twochambers} would
give an integral determinant-one matrix sending it into $P_0$ or $P_1$.
For a source replaced by $G_jp$ in target $j$, the corresponding matrix
is $HG_j^{-1}$. The complete search for that target would retain it,
contradicting its empty result. The positive realizations therefore
give the asserted equality. No arithmetic-normalizer property of the
reflection and no independent per-coordinate chamber choice is used.
\end{proof}

